% A6 manual-test fixture — item 7 (a cut-off answer says what it lost).
% The Method and Experiments sections are deliberately empty scaffolds: the test
% asks Lemmaria to draft both in ONE edit, which is what pushes a single turn
% past the 8192-token per-turn output ceiling.
\documentclass[11pt]{article}

\usepackage[utf8]{inputenc}
\usepackage{amsmath}
\usepackage{amssymb}
\usepackage{booktabs}
\usepackage{graphicx}
\usepackage{natbib}
\usepackage{hyperref}

\title{Corridor Attention: Bounded-Bandwidth Routing for Long-Context
Sequence Models}
\author{A.~Fixture \and B.~Fixture}
\date{}

\begin{document}
\maketitle

\begin{abstract}
Transformer language models spend a quadratic share of their compute on
attention, which puts a hard practical ceiling on the context lengths that are
affordable to train. We introduce \emph{Corridor Attention}, a routing scheme
that gives every query token a fixed bandwidth of $k$ key positions drawn from
a learned, block-level corridor rather than from the full prefix. Corridor
Attention keeps a dense local window for short-range syntax, adds a small set
of globally visible summary tokens for discourse-level dependencies, and routes
the remaining bandwidth through a differentiable block scorer that is trained
jointly with the rest of the network. The result is an attention layer whose
cost is linear in sequence length with a constant that the practitioner sets
directly. We describe the method, characterise its cost, and evaluate it on
long-context language modelling and on retrieval-style probes.
\end{abstract}

\section{Introduction}\label{sec:intro}

Self-attention \citep{vaswani2017attention} computes an interaction between
every pair of positions in a sequence. That property is the reason
transformers model long-range structure well, and it is also the reason their
cost grows as $\Theta(n^2)$ in the sequence length $n$. For the context
lengths that are now routinely requested of production models --- tens of
thousands of tokens of source code, a whole scientific article and its
appendices, a multi-day conversation --- the attention term stops being a
constant factor and starts being the budget.

The literature has answered this in two broadly different ways. One family
\emph{restricts} the interaction pattern: sliding windows, dilated windows,
strided or block-sparse masks, and hand-designed combinations of the three
\citep{child2019sparse,beltagy2020longformer,zaheer2020bigbird}. These methods
are cheap and predictable, but the pattern is fixed before the model sees any
data, so a dependency the designer did not anticipate simply cannot be
expressed. The other family \emph{approximates} the softmax kernel, replacing
the exact attention matrix with a low-rank or feature-map surrogate
\citep{wang2020linformer,katharopoulos2020linear,choromanski2021performer}.
These methods keep a formally global receptive field, but the approximation
error is not under the practitioner's control and it interacts badly with the
sharply peaked attention distributions that emerge late in training.

We take a third position, which we think is the pragmatic one: keep attention
\emph{exact}, and make the \emph{choice of what to attend to} learned. Each
query is allocated a fixed bandwidth --- a corridor --- of $k$ key positions.
A lightweight scorer, operating on block summaries rather than on individual
positions, decides which blocks the corridor passes through. Because the
bandwidth is a hyperparameter rather than an emergent property, the cost of the
layer is known before training starts, and because the routing is learned, the
pattern can differ per head, per layer, and per input.

Three design commitments follow from this position, and they are what
distinguishes Corridor Attention from the closest prior work.

\paragraph{Bandwidth is declared, not discovered.} The practitioner sets $k$.
The model may not exceed it, and the training objective contains no term that
trades accuracy against cost --- there is nothing to trade, because the cost is
fixed. This makes capacity planning a matter of arithmetic rather than of
empirical scaling curves.

\paragraph{Routing is scored on blocks, decided per query.} Scoring every
candidate key position would reintroduce the quadratic term we are trying to
remove. Corridor Attention scores $\lceil n / b \rceil$ block summaries
instead, at a cost of $\Theta(n^2 / b^2)$ for the scorer, which for the block
sizes we use is a small fraction of the total. The selection itself is
per-query, so two queries in the same block may route differently.

\paragraph{The local window is not negotiable.} Ablating the dense local
window degrades perplexity far more than ablating any equivalent amount of
routed bandwidth. We therefore treat the window as part of the architecture
rather than as one option among several, and we spend the learned bandwidth
only on what the window cannot reach.

Our contributions are as follows. (i) We formalise bounded-bandwidth routed
attention and give a block-scored, per-query instantiation of it that trains
end to end without an auxiliary load-balancing loss. (ii) We give a cost model
for the layer and validate it against measured wall-clock time, showing that
the predicted and observed crossover points against dense attention agree
closely. (iii) We evaluate on long-context language modelling and on synthetic
retrieval probes, and we report the ablations that identify which of the three
components each benefit comes from.

\section{Background and Notation}\label{sec:background}

Let $X \in \mathbb{R}^{n \times d}$ be the input to an attention layer, with
$n$ positions and model width $d$. A single head projects $X$ into queries,
keys and values, $Q = XW_Q$, $K = XW_K$, $V = XW_V$, with
$W_Q, W_K, W_V \in \mathbb{R}^{d \times d_h}$, and computes
\begin{equation}\label{eq:dense}
  \mathrm{Attn}(Q, K, V) = \mathrm{softmax}\!\left(
    \frac{QK^{\top}}{\sqrt{d_h}} + M \right) V,
\end{equation}
where $M \in \{0, -\infty\}^{n \times n}$ is the causal mask. Equation
\eqref{eq:dense} requires $\Theta(n^2 d_h)$ multiply--accumulates and, in a
naive implementation, $\Theta(n^2)$ memory; fused kernels
\citep{dao2022flashattention} remove the memory term but not the compute term.

A \emph{sparse} attention layer replaces $M$ with a mask $M_S$ whose finite
entries are indexed by a set $S \subseteq \{1,\dots,n\}^2$ of permitted
query--key pairs. We call $|S_i| = |\{j : (i,j) \in S\}|$ the
\emph{bandwidth} of query $i$. Fixed-pattern methods choose $S$ in advance.
Corridor Attention chooses $S$ as a function of the input, subject to
$|S_i| \le k$ for a constant $k$ set by the practitioner.

We write $b$ for the block size used by the router, $m = \lceil n/b \rceil$ for
the number of blocks, $w$ for the width of the dense local window, and $g$ for
the number of global summary tokens. Throughout, \emph{bandwidth} refers to
$k = w + g + r$, where $r$ is the routed bandwidth.

\section{Related Work}\label{sec:related}

\paragraph{Fixed sparse patterns.} Sparse Transformer
\citep{child2019sparse} introduced strided and fixed factorised patterns.
Longformer \citep{beltagy2020longformer} combined a sliding window with
task-specific global tokens, and BigBird \citep{zaheer2020bigbird} added
random attention and gave the resulting pattern a theoretical justification.
All three fix the pattern before training. Our local window and global summary
tokens are drawn directly from this line of work; the routed component is what
we add to it.

\paragraph{Kernel and low-rank approximations.} Linformer
\citep{wang2020linformer} projects keys and values to a fixed rank. Linear
attention \citep{katharopoulos2020linear} and Performer
\citep{choromanski2021performer} replace the softmax with a feature map,
recovering linear cost and a recurrent formulation. These methods approximate
the attention matrix; Corridor Attention computes an exact softmax over a
restricted support, which we prefer because the error is then structural and
inspectable rather than distributional.

\paragraph{Learned and content-based routing.} Routing Transformer
\citep{roy2021efficient} clusters queries and keys with online $k$-means and
attends within clusters. Reformer \citep{kitaev2020reformer} uses
locality-sensitive hashing to the same end. Both make the \emph{grouping}
content-dependent but leave each query's bandwidth as a consequence of cluster
occupancy, which is why both need balancing machinery. Corridor Attention fixes
the bandwidth per query and lets the routing vary, which we argue is the more
useful of the two invariants to hold fixed.

\paragraph{Retrieval-augmented and memory models.} A separate line of work
keeps long context outside the attention layer entirely, in a retrieval index
or a compressed memory \citep{rae2020compressive,borgeaud2022improving}. These
are complementary to our method rather than alternatives to it: Corridor
Attention is what we would use \emph{inside} such a model to read a long
retrieved passage.

\section{Method}\label{sec:method}

% TODO(draft): Corridor Attention, in full. Cover, as \subsection blocks:
%   1. Overview and notation recap.
%   2. Block summaries: how a block summary is computed and why mean-pooling
%      the keys is sufficient.
%   3. The corridor scorer: architecture, inputs, and the top-r selection.
%   4. Straight-through training of the discrete selection; why no auxiliary
%      load-balancing loss is needed.
%   5. Composing the three components (local window, global summaries, routed
%      corridor) into one exact softmax over the union of supports.
%   6. Cost model: FLOPs and memory as functions of n, b, w, g, r, and the
%      predicted crossover against dense attention.
%   7. Implementation notes: the gather/scatter kernel, and what makes the
%      routed component efficient on real hardware.

\section{Experiments}\label{sec:experiments}

% TODO(draft): the full experimental section. Cover, as \subsection blocks:
%   1. Datasets and preprocessing (long-form text, code, and the synthetic
%      retrieval probes).
%   2. Baselines: dense attention, sliding window, Longformer-style,
%      BigBird-style, Linformer, Performer, Routing Transformer.
%   3. Training setup: model sizes, optimiser, schedule, hardware.
%   4. Main results: perplexity vs. context length, and vs. wall-clock cost.
%   5. Retrieval probes: needle-in-a-haystack style accuracy by distance.
%   6. Ablations: remove the local window; remove global summaries; freeze the
%      router at initialisation; vary b, w, g and r independently.
%   7. Cost-model validation: predicted vs. measured throughput.
%   8. Limitations and negative results.

\section{Discussion}\label{sec:discussion}

Fixing bandwidth and learning the pattern is, in our reading, the trade that
practitioners actually want. A fixed pattern gives a predictable bill and an
unpredictable failure mode; a learned pattern with a free bandwidth gives the
reverse. Corridor Attention makes the bill predictable and moves the learned
degree of freedom to where it costs nothing at inference time.

Two limitations are worth stating plainly. First, the router is trained with a
straight-through estimator, and although we observe stable training across all
the scales we tried, we have no guarantee that this continues to hold at
larger widths. Second, a fixed bandwidth is by construction the wrong choice
for inputs whose genuine dependency structure is heavy-tailed; for those, a
method that can spend more on a few queries would be preferable, at the cost of
the predictability we set out to obtain.

\section{Conclusion}\label{sec:conclusion}

We presented Corridor Attention, an exact-softmax attention layer with a
declared per-query bandwidth and a learned, block-scored routing of that
bandwidth. It keeps the cost of long-context attention linear with a constant
the practitioner chooses, and it recovers most of the quality of dense
attention on the long-context benchmarks we tested.

\bibliographystyle{plainnat}
\bibliography{refs}

\end{document}
